\SourcePage{409}{412}
\section{$\Lambda$-doubling}

The twofold degeneracy of terms with $\Lambda\ne0$ (see \S~78) is actually approximate. It holds only insofar as the influence of molecular rotation on the electronic state (and higher-order spin--orbit effects) is neglected, as throughout the preceding theory. Allowing for the interaction of the electronic state with rotation splits a term with $\Lambda\ne0$ into two closely spaced levels. This phenomenon is called $\Lambda$-doubling (E.~Hill, J.~van Vleck, R.~Kronig, 1928).

\SourcePage{410}{413}
We again begin the quantitative treatment with singlet terms ($S=0$). In \S~82 the rotational-level energies were calculated to first order in perturbation theory by finding diagonal matrix elements (the mean value) of the operator
\[
B(r)(\widehat{\mathbf K}-\widehat{\mathbf L})^2.
\]
For the subsequent orders we must consider elements of this operator nondiagonal in $\Lambda$. The operators $\widehat{\mathbf K}^2$ and $\widehat{\mathbf L}^2$ are diagonal in $\Lambda$, so only $2B\widehat{\mathbf K}\widehat{\mathbf L}$ need be considered.

The matrix elements of $\widehat{\mathbf K}\widehat{\mathbf L}$ are conveniently calculated with (29.12), setting $\mathbf A=\mathbf K$, $\mathbf B=\mathbf L$; $K,M_K$ play the roles of $L,M$, while $n,\Lambda$ replaces $n$, where $n$ denotes the set of quantum numbers (excluding $\Lambda$) that specifies the electronic term. Since the matrix of the conserved vector $\mathbf K$ is diagonal in $n,\Lambda$, while the matrix of $\mathbf L$ has nondiagonal elements only for transitions changing $\Lambda$ by one (compare the discussion of an arbitrary vector $\mathbf A$ in \S~87), formulas (87.3) give
\begin{equation}
\langle n'\Lambda KM_K|\mathbf{KL}|n,\Lambda-1,KM_K\rangle
=\frac12\langle n'\Lambda|L_\xi+iL_\eta|n,\Lambda-1\rangle
\sqrt{(K+\Lambda)(K+1-\Lambda)}.
\tag{88.1}
\end{equation}
There are no matrix elements corresponding to a larger change in $\Lambda$.

The perturbing action of the elements for $\Lambda\to\Lambda-1$ can produce an energy difference between the states with $\pm\Lambda$ only in order $2\Lambda$ of perturbation theory. The effect is consequently proportional to $B^{2\Lambda}$, that is, to $(m/M)^{2\Lambda}$ ($M$ is the nuclear mass and $m$ the electron mass). For $\Lambda>1$ this quantity is too small to be of interest. Thus $\Lambda$-doubling is significant only for $\Pi$ terms ($\Lambda=1$), which are considered below.

For $\Lambda=1$ we must go to second order. The corrections to the energy eigenvalues may be found from the general formula (38.10). The denominators of its summands contain energy differences of the form $E_{n,\Lambda,K}-E_{n',\Lambda-1,K}$. The terms involving $K$ cancel in these differences, because at a fixed internuclear distance $r$ the rotational energy has the same value $B(r)K(K+1)$ for every term.

\SourcePage{411}{414}
The dependence of the required splitting $\Delta E$ on $K$ is therefore determined entirely by the squares of the matrix elements in the numerators. These include squares of the elements for transitions changing $\Lambda$ from 1 to 0 and from 0 to $-1$; by (88.1) both have the same dependence on $K$. The splitting of a ${}^1\Pi$ term is consequently
\begin{equation}
\Delta E=\mathop{\rm const}\nolimits\,K(K+1),
\tag{88.2}
\end{equation}
where, in order of magnitude, $\mathop{\rm const}\nolimits\sim B^2/\varepsilon$, and $\varepsilon$ is of the order of the separations between neighboring electronic terms.

We turn to terms with nonzero spin (${}^2\Pi$ and ${}^3\Pi$ terms; higher values of $S$ scarcely occur). If the term belongs to case $b$, the multiplet splitting has no effect at all on the $\Lambda$-doubling of the rotational levels, which is still given by (88.2).

In case $a$, by contrast, spin has an essential influence. Each electronic term is characterized not only by $\Lambda$ but also by $\Omega$. Merely replacing $\Lambda$ by $-\Lambda$ changes $\Omega=\Lambda+\Sigma$, and therefore gives an entirely different term. The mutually degenerate states are those with $\Lambda,\Omega$ and $-\Lambda,-\Omega$. This degeneracy can be removed not only by the interaction between orbital angular momentum and molecular rotation considered above, but also by the spin--orbit interaction. Conservation of the projection $\Omega$ of the total angular momentum on the molecular axis is an exact conservation law (for fixed nuclei), and hence cannot be violated by the spin--orbit interaction. The latter can, however, change $\Lambda$ and $\Sigma$ simultaneously, having matrix elements for the corresponding transitions, in such a way that $\Omega$ remains unchanged. Acting alone, or together with the orbit--rotation interaction (which changes $\Lambda$ without changing $\Sigma$), this effect can produce $\Lambda$-doubling.

First consider ${}^2\Pi$ terms. For a ${}^2\Pi_{1/2}$ term ($\Lambda=1$, $\Sigma=-1/2$, $\Omega=1/2$), the splitting arises when both the spin--orbit and orbit--rotation interactions are included, each to first order. The former produces the transition $\Lambda=1$, $\Sigma=-1/2\to\Lambda=0$, $\Sigma=1/2$; the latter then takes the state $\Lambda=0$, $\Sigma=1/2$ into $\Lambda=-1$, $\Sigma=1/2$, which differs from the initial state by reversal of the signs of $\Lambda$ and $\Omega$. The spin--orbit matrix elements do not depend on the rotational quantum number $J$, while the dependence of the orbit--rotation matrix elements is given by (88.1), with $K,\Lambda$ under the radical replaced by $J,\Omega$. Thus

\SourcePage{412}{415}
the $\Lambda$-doubling of the ${}^2\Pi_{1/2}$ term is
\begin{equation}
\Delta E_{1/2}=\mathop{\rm const}\nolimits\,(J+1/2),
\tag{88.3}
\end{equation}
where $\mathop{\rm const}\nolimits\sim AB/\varepsilon$. For a ${}^2\Pi_{3/2}$ term, however, the splitting can arise only in higher orders, so in practice $\Delta E_{3/2}=0$.

Finally consider ${}^3\Pi$ terms. For a ${}^3\Pi_0$ term ($\Lambda=1$, $\Sigma=-1$), the splitting is produced by the spin--orbit interaction in second order (through the transitions $\Lambda=1$, $\Sigma=-1\to\Lambda=0$, $\Sigma=0\to\Lambda=-1$, $\Sigma=1$). Hence in this case the $\Lambda$-doubling is entirely independent of $J$:
\begin{equation}
\Delta E_0=\mathop{\rm const}\nolimits,
\tag{88.4}
\end{equation}
where $\mathop{\rm const}\nolimits\sim A^2/\varepsilon$. For a ${}^3\Pi_1$ term, $\Sigma=0$, so spin does not affect the splitting and we recover a formula of the form (88.2), with $K$ replaced by $J$:
\begin{equation}
\Delta E_1=\mathop{\rm const}\nolimits\,J(J+1).
\tag{88.5}
\end{equation}
For a ${}^3\Pi_2$ term still higher orders are required, and one may take $\Delta E_2=0$.

Of the two levels produced by $\Lambda$-doubling, one is always positive and the other negative; this was already discussed in \S~6. An analysis of the molecular wave functions determines the pattern in which positive and negative levels alternate. We give only the result of that analysis here.\footnote{See the article by Wigner and Witmer cited on p.~376.} If for some $J$ the positive level lies below the negative one, their order is reversed in the doublet with $J+1$: the positive level lies above the negative one, and so on. The ordering alternates with successive values of the total angular momentum (for case $a$ terms; in case $b$ the same statement holds for successive values of $K$).

\ProblemHeading Determine the $\Lambda$-splitting of a ${}^1\Delta$ term.

\SolutionHeading Here the effect first appears in fourth-order perturbation theory. Its dependence on $K$ is determined by products of four matrix elements (88.1), for transitions changing $\Lambda$ as $2\to1$, $1\to0$, $0\to-1$, $-1\to-2$. This gives
\[
\Delta E=\mathop{\rm const}\nolimits\,(K-1)K(K+1)(K+2),
\]
where $\mathop{\rm const}\nolimits\sim B^4/\varepsilon^3$.

\SourcePage{413}{416}
